\documentclass[11pt]{article}
\usepackage[a4paper,margin=27mm]{geometry}
\usepackage{amsmath,amssymb,mathtools}
\usepackage{fontspec,unicode-math}
\setmainfont{Latin Modern Roman}
\setmathfont{Latin Modern Math}
\usepackage{xurl}
\usepackage[hidelinks]{hyperref}
\newcommand{\sref}[2]{\hyperlink{source:#1:#2}{[S#2]}}
\newcommand{\cataloguescope}[1]{\paragraph{Mathematical question.}#1}
\hypersetup{pdftitle={UnsolvedMath Problem 9000004}}
\begin{document}
% UnsolvedMath ID 9000004; source catalogue code AMR-089-0004
\setcounter{section}{9000003}
\section{Alternative and arithmetic proofs of the pi identities}\label{9000004}
{\small\sffamily UnsolvedMath ID 9000004 · Graph Theory}
\subsection{Definitions and mathematical statement}

\paragraph{Shared notation.}$\mathbb N=\{1,2,\ldots\}$, and $\mathbb Z,\mathbb Q,\mathbb R,\mathbb C$ denote the integers, rationals, reals and complex numbers. Any other conventions are specified locally.


Let $\mathcal A=\{(a,b,c,d)\in\mathbb Z_{\ge0}^{4}:ad-bc=1\}$ and put
\[
f(a,b,c,d)=\sqrt{a^2+b^2}+\sqrt{c^2+d^2}-\sqrt{(a+c)^2+(b+d)^2}.
\]
The two identities in question are
\[
\sum_{\mathcal A}f(a,b,c,d)=2,\qquad
\sum_{\mathcal A}f(a,b,c,d)^2=2-\frac\pi2.
\]
The comparison identity is the Basel formula $\sum_{n=1}^{\infty}n^{-2}=\pi^2/6$. In the source its methods include infinite-product coefficient comparisons and complex residues. The Gaussian integers are $\mathbb Z[i]=\{m+ni:m,n\in\mathbb Z\}$, where $i^2=-1$. The residue of a meromorphic function $g$ at $z_0$ is the coefficient of $(z-z_0)^{-1}$ in its Laurent series.
\paragraph{Statement recorded in the supplied source.}
Find other proofs of both displayed lattice-sum identities using methods of proof of the Basel formula. Can the individual summand $f(a,b,c,d)$ be interpreted as a residue of an appropriate complex function at $(a+b)+(c+d)i$, or otherwise related arithmetically to Gaussian integers? \sref{9000004}{2}
\subsection{Short English statement}
Give Basel-style proofs of both lattice identities, and find a residue or Gaussian-integer interpretation of their summand.
\subsection{Sources}
\begin{description}
\item[\hypertarget{source:9000004:1}{S1}] ULAM.AI and the UnsolvedMath Contributors, UnsolvedMath: A Curated Collection of Open Mathematics Problems, record 9000004 (AMR-089-0004). CC BY 4.0. \url{https://huggingface.co/datasets/ulamai/UnsolvedMath}.
\item[\hypertarget{source:9000004:2}{S2}] Nikita Kalinin and Mikhail Shkolnikov, \emph{The number $\pi$ and summation by $SL(2,\mathbb Z)$}, Arnold Mathematical Journal (2018) \S4, Questions, pp. 5--6. \url{https://amj.math.stonybrook.edu/html-articles/Files-2015-2024/17-75/index.html}.
\end{description}
\label{end:9000004}
\end{document}
